\subsection*{Transfer to classroom questions}
We replaced a training distribution dominated by explicit grammatical terms with a mixture that included conversational questions about actions and their objects. The passage images, answer images and model size stayed fixed. On 539 previously held passages, the new reader answered both English and both Chinese classroom questions correctly on 98 passages. The previous reader and the exchanged-answer control each scored zero. This gain concealed an asymmetry. The Chinese pair improved from 4 to 240 complete passages, whereas the English pair declined from 135 to 107. On the original eight order-sensitive requests the score fell from 230 to 227.

The stricter test used 29 previously unevaluated excerpts from the \textit{Zhuangzi} and \textit{Han Feizi} (Fig.~\ref{fig:everyday-gloss}). Only one passed all four questions. For the passage about returning a pearl, the model drew \textit{return} or \textit{pearl} as requested. It located both words on six passages across both languages. None of the 14 passages with a fronted object passed this location test. Even when supplied with the correct locations, the writer could give both registered English readings on only five of the 29 passages. For one printed object pronoun it drew the verb \textit{go}. Some misses reflect narrow reference lists: the writer produced \textit{amass} where the registered list allowed \textit{gather}, \textit{collect} and \textit{accumulate}. We retained that miss. The results support limited transfer of instructions while exposing separate errors in word order, lexical sense and spelling.
